\section{A Gaussian rate function need not determine the conditioned posterior}
\label{sec:non-gaussian-tie}

Moderate deviations do justify a Gaussian rate function throughout
$0<\gamma<1$. They do not, however, justify replacing the actual
posterior by a conditioned Gaussian in total variation throughout this
interval. The distinction already occurs with three input bits and one
hidden layer. It affects the criterion that the correct label must
have posterior probability greater than $2/3$.

We first give the estimate that makes the counterexample rigorous.
All vectors in the following lemma have a fixed dimension $m$.

\begin{lemma}[Annealed Gaussian corner probabilities]
\label{lem:annealed-corner}
For every $N$, let $H_{N,1},\ldots,H_{N,N}$ be independent copies of a
random vector $H_N\in[-1,1]^m$. Suppose
\[
 K_N=\mathbb E H_NH_N^T\longrightarrow K>0,
 \qquad H_N\ \Longrightarrow\ H.
\]
Set $Q_N=N^{-1}\sum_i H_{N,i}H_{N,i}^T$, and, conditionally on the
features, let $Z_N\sim\mathcal N(0,t_N^{-1}Q_N)$, where
$t_N\to\infty$ and $t_N/N\to0$.
Fix $a>0$ and $g\in\{-1,1\}^m$. Assume
\[
 g_i(K^{-1}g)_i>0\quad\hbox{for every }i.
\]
Writing $v=K^{-1}g$, define
\[
 E_{N,g}=\frac{a^2}{2}g^TK_N^{-1}g,
 \qquad
 C_g=(2\pi)^{-m/2}(\det K)^{-1/2}
             \prod_{i=1}^m\bigl[a g_i v_i\bigr]^{-1}.
\]
Then
\begin{align}
 \log\Pr\{g_iZ_{N,i}>a\ \forall i\}
 ={}&-t_NE_{N,g}-\frac m2\log t_N+\log C_g\notag\\
 &+\frac{a^4t_N^2}{8N}
       \operatorname{Var}\bigl((v^TH)^2\bigr)
       +o\bigl(1+t_N^2/N\bigr).
 \label{eq:annealed-corner}
\end{align}
In particular, the leading correction to the deterministic Gaussian
corner probability can grow without bound even though $t_N=o(N)$.
\end{lemma}

\begin{proof}
Write $D_g=\operatorname{diag}(g)$ and
$E_g(Q)=a^2g^TQ^{-1}g/2$. In a fixed neighborhood of $K$, the
matrices $Q$ are positive definite and all entries of
$\beta_g(Q)=aD_gQ^{-1}g$ are bounded below by a positive constant.
The change of variables $z=ag+D_gu/t_N$ in the Gaussian density gives
\begin{align*}
 \Pr_{Q}\{g_iZ_i>a\ \forall i\}
 ={}&(2\pi)^{-m/2}(\det Q)^{-1/2}t_N^{-m/2}e^{-t_NE_g(Q)}\\
 &\quad\times\int_{[0,\infty)^m}
  \exp\left\{-\beta_g(Q)^Tu
       -\frac{u^TD_gQ^{-1}D_gu}{2t_N}\right\}\,du.
\end{align*}
The integral equals $\prod_i\beta_g(Q)_i^{-1}[1+O(t_N^{-1})]$
uniformly in this neighborhood: use
$0\le1-e^{-r}\le r$ and the finite second moments of the product
exponential density. Denote the resulting smooth positive prefactor
by $C_g(Q)$.

Put $v_N=K_N^{-1}g$ and $\Delta_N=Q_N-K_N$. Matrix inversion has the
uniform expansion
\[
 -t_NE_g(Q_N)
 =-t_NE_{N,g}+t_N\ell_N(\Delta_N)
                  +O(t_N\|\Delta_N\|^2),
 \qquad
 \ell_N(\Delta)=\frac{a^2}{2}v_N^T\Delta v_N.
\]
The linear statistic is the average of the bounded centered variables
\[
 Y_{N,i}=\frac{a^2}{2}
   \left[(v_N^TH_{N,i})^2-
                   \mathbb E(v_N^TH_N)^2\right].
\]
Consequently, with $\theta_N=t_N/N$,
\begin{align*}
 \log\mathbb E e^{t_N\ell_N(\Delta_N)}
 &=N\log\mathbb E e^{\theta_NY_{N,1}}\\
 &=\frac{t_N^2}{2N}\operatorname{Var}(Y_{N,1})
                              +O(t_N^3/N^2).
\end{align*}
All bounds here are uniform, since the $Y_{N,i}$ are uniformly bounded.

For completeness, the quadratic remainder is negligible at the
claimed precision, even when $t_N^2/N$ diverges. Tilt the independent
feature rows by the density proportional to
$\exp(\theta_NY_{N,i})$. Under this product law their mean covariance
differs from $K_N$ by $O(\theta_N)$, and Hoeffding's inequality still
gives, for a fixed number of matrix entries,
\[
 \Pr_{\theta_N}\{\|\Delta_N-
       \mathbb E_{\theta_N}\Delta_N\|>r\}
       \le C e^{-cNr^2}.
\]
Integrating this bound, and using $t_N/N\to0$, yields
\[
 \log\mathbb E_{\theta_N}
   \exp\{C t_N\|\Delta_N\|^2+C\|\Delta_N\|\}
 =O\left(\frac{t_N^3}{N^2}+\frac{t_N}{N}
                                      +N^{-1/2}\right).
\]
A matching lower estimate for the negative of the expression follows
from Jensen's inequality. These bounds control both the inversion
remainder and the smooth prefactor $C_g(Q_N)/C_g(K_N)$.

The complement of the fixed covariance neighborhood has probability
$O(e^{-cN})$, by entrywise Hoeffding. Its contribution to the original
orthant probability is at most this quantity. Inside the neighborhood
the calculation gives a probability $\exp[-O(t_N)-O(\log t_N)]$;
hence the omitted contribution is negligible. The same truncation
under the tilted law is harmless, since the tilt is $o(N)$ and the
features are bounded. We have thus obtained the expansion with
$C_g(K_N)$, $v_N$, and $H_N$, and an error
\[
 O\left(t_N^{-1}+N^{-1/2}+t_N/N+t_N^3/N^2\right)
                  =o(1+t_N^2/N).
\]
Finally $C_g(K_N)\to C_g$, and boundedness and weak convergence give
$\operatorname{Var}((v_N^TH_N)^2)\to
\operatorname{Var}((v^TH)^2)$. This proves the displayed claim.
\end{proof}

\begin{theorem}[An actual-network counterexample to Gaussian posterior equivalence]
\label{thm:gaussian-posterior-counterexample}
Use the original architecture with $d=3$, one hidden layer ($L=2$),
activation $\tanh$, and any fixed margin $a>0$. Order the cube inputs
lexicographically, starting with $x_0=(-1,-1,-1)$. Define
\[
 g=(-1,1,1,-1,-1,-1,-1,1),\qquad
 h=(1,1,1,-1,-1,-1,-1,1).
\]
Observe the seven common labels at every input other than $x_0$.
Let $p_N$ be the actual posterior probability of $g$, after the
all-input-margin restriction and these observations. Let
$\widetilde p_N$ be the corresponding probability for the deterministic
Gaussian output vector $\mathcal N(0,N^{-\gamma}K_N)$, where $K_N$
is the exact mean feature covariance at width $N$. Then
\[
 \widetilde p_N\longrightarrow \frac{189}{314}
                  \qquad (0<\gamma<1),
\]
whereas
\[
 p_N\longrightarrow
 \begin{cases}
  189/314,&0<\gamma<1/2,\\[2pt]
  \displaystyle\left[1+\frac{125}{189}
              e^{-8586a^4/625}\right]^{-1},&\gamma=1/2,\\[6pt]
  1,&1/2<\gamma<1.
 \end{cases}
\]
The same Gaussian limit holds if $K_N$ is replaced by the limiting
arcsine kernel $K$. Thus, for $1/2\le\gamma<1$, the actual Boolean
posterior need not converge in total variation to either of these
conditioned Gaussian posteriors.

For the original margin $a=1/2$ and target $g$, the Gaussian learner's
correct-label probability at the unseen input tends to
$189/314\simeq0.602$, whereas the actual learner's probability tends
to approximately $0.781$ at $\gamma=1/2$ and to one for
$\gamma>1/2$. The proposed $2/3$ confidence criterion therefore
separates their predictions.
\end{theorem}

\begin{proof}
Let $A$ be the $8\times4$ matrix with rows $(1,x)$, in the specified
order; then $A^TA=8I_4$. A feature row has distribution
\[
 H_N=\tanh\left(\frac{N^{(1-\gamma)/2}}{\sqrt3}AW\right),
                 \qquad W\sim\mathcal N(0,I_4).
\]
For $\gamma<1$ it converges almost surely to $H=\operatorname{sign}(AW)$.
Gaussian symmetry and oddness show, exactly at every finite width,
\[
 K_N=(q_N-2c_N)I_8+\frac{c_N}{2}AA^T,
\]
where $q_N$ is the diagonal feature covariance and $c_N$ is the
covariance between cube neighbors. The other off-diagonal entries are
zero at Hamming distance two and $-c_N$ at distance three. Moreover
$q_N\to1$ and $c_N\to1/3$, so
\[
 K=\frac13I_8+\frac16AA^T,\qquad
 K^{-1}=3I_8-\frac3{10}AA^T.
\]
The two functions satisfy
\[
 A^Tg=(-2,-2,2,2)^T,\qquad A^Th=(0,-4,0,0)^T.
\]
Their squared lengths are both eight and the displayed projection
vectors both have squared length sixteen. Therefore
$g^TK_N^{-1}g=h^TK_N^{-1}h$ exactly for all sufficiently large $N$;
the finite-gain correction cannot break the energy tie.

Direct multiplication gives
\begin{align*}
 v:=K^{-1}g&=\frac35(-3,5,5,-7,-1,-3,-3,5),\\
 w:=K^{-1}h&=\frac35(3,3,3,-7,-3,-3,-3,7).
\end{align*}
All signed gradient coordinates are strictly positive, and
$g^Tv=h^Tw=96/5$. Both orthant minimizers are therefore their corners
$ag$ and $ah$. The Gaussian corner prefactors have ratio
\[
 \frac{C_g}{C_h}
 =\frac{\prod_i h_i w_i}{\prod_i g_i v_i}
 =\frac{3^6\,7^2}{3^3\,5^3\,7}=\frac{189}{125}.
\]
The uniform Gaussian corner estimate in
Lemma~\ref{lem:annealed-corner}, before averaging over features, proves
$\widetilde p_N\to189/314$.

It remains to compute a fourth moment of the saturated feature row.
Here this calculation is exact and finite. Signed permutations of the
four coordinates of $W$ give three orbits of possible threshold
vectors $H$, of sizes $8,64,32$. To check completeness, first sort the
absolute coordinates as $r_1>r_2>r_3>r_4>0$. Only the signs of
$r_1-r_2-r_3-r_4$ and $r_1-r_2-r_3+r_4$ are undetermined; their
three feasible sign combinations give precisely these three orbits.
The exceptional equalities have probability zero.

Let $P=AA^T/8$. The following table records finite averages within
each orbit. Representatives are coefficient vectors $W$, not feature
vectors. The column $n_+$ counts orbit elements with coordinates
$H_1,H_2,H_4,H_7$ all positive, using zero-based cube indices.
\[
\begin{array}{c|r|r|r|r|r}
 \text{representative }W&\text{size}&\|PH\|^2&n_+
 &\text{mean }(v^TH)^4&\text{mean }(w^TH)^4\\ \hline
 (10,3,2,1)&8&8&1&1296/625&5184/625\\
 (5,4,2,3/2)&64&6&4&964224/625&1174176/625\\
 (4,7/2,3,1)&32&6&0&1174176/625&754272/625
\end{array}
\]
Each entry can be checked by applying the $2^4 4!$ signed
permutations to the representative, taking $\operatorname{sign}(AW)$,
removing duplicates, and summing the indicated integer fourth powers.
Gaussian invariance makes the distribution uniform within each orbit.
To determine the three orbit probabilities $p_1,p_2,p_3$, note that
\[
 8p_1+6p_2+6p_3
       =\mathbb E\|PH\|^2=\operatorname{tr}(PK)=20/3.
\]
Thus $p_1=1/3$. The four rows of $A$ with indices $1,2,4,7$ are
orthogonal, so their Gaussian projections have independent signs.
Using the $n_+$ column,
\[
 \frac1{16}=\Pr\{H_1=H_2=H_4=H_7=1\}
            =\frac{p_1}{8}+\frac{p_2}{16}.
\]
It follows that $p_2=p_3=1/3$. The table consequently gives
\[
 \mathbb E(v^TH)^4=\frac{713232}{625},\qquad
 \mathbb E(w^TH)^4=\frac{644544}{625}.
\]
The second moments are equal, both being $96/5$. Hence the difference
of the two variances of squared projections is $68688/625$.

Apply Lemma~\ref{lem:annealed-corner} with $t_N=N^\gamma$.
Since the two energies agree exactly, their prior log odds satisfy
\[
 \log\frac{\Pr(A_g)}{\Pr(A_h)}
 =\log\frac{189}{125}
    +\frac{8586a^4}{625}N^{2\gamma-1}
    +o(1+N^{2\gamma-1}).
\]
The seven observed labels leave exactly these two Boolean functions
compatible, so these are also their posterior odds. The three stated
limits follow directly.
\end{proof}

\begin{remark}[What this counterexample does and does not show]
The leading kernel margin energy remains correct at every
$0<\gamma<1$: both functions in the example have the same energy, and
the actual posterior selects one of these tied minimizers. The
counterexample disproves full conditional-Gaussian equivalence, even
with the exact finite-width mean covariance. It does not by itself
establish a circuit prior or a general sample-complexity advantage.
It shows that the leading kernel energy alone need not determine the
posterior's confidence, and that fourth-order feature statistics can
remain visible after rare-event conditioning.
\end{remark}
